\documentclass[12pt]{article}
\usepackage{amsmath,amssymb}
\newcommand{\Dr}{\operatorname{Dr}}
\newcommand{\Trop}{\operatorname{Trop}}
\newcommand{\rk}{\operatorname{rk}}
\newcommand{\cl}{\operatorname{cl}}
\newcommand{\cL}{\mathcal L}
\begin{document}
\section*{Research notes}

Problem: Given a matroid $M$ and the relative Dressian $\Dr(M)$ of
tropical linear subspaces contained in $\Trop(M)$, suppose the lattice
of flats $\cL(M)$ has an order-reversing involution
$F\mapsto F^\perp$.  Need an order-reversing involution on $\Dr(M)$
extending $F\mapsto F^\perp$ under the embedding
\[
F\longmapsto \Trop(M/F\oplus U_{0,F}).
\]

\subsection*{Conventions}

Use projective tropical linear spaces as in Brandt--Eur--Zhang.  A
rank $k$ element of $\Dr(M)$ is represented by a rank $k$ valuated
matroid $\nu$ that is a valuated quotient of the trivial valuated
matroid $\mu_M$ of rank $r=\rk M$.  Inclusion of projective tropical
linear spaces corresponds to quotient order:
\[
\Trop(\theta)\subseteq\Trop(\nu)
\Longleftrightarrow
\theta \text{ is a valuated quotient of } \nu.
\]
Valuated matroids are projective: adding a common scalar to all
Plucker coordinates does not change the tropical linear space.

For valuated matroids $p,q$ of ranks $a<b$ on the same ground set, the
incidence relation used throughout is
\[
\mathcal I_{a,b}(p,q;S,T):\quad
\min_{i\in T\setminus S}\{p(S\cup i)+q(T\setminus i)\}
\text{ is attained at least twice},
\]
where $|S|=a-1$, $|T|=b+1$, and an all-$\infty$ relation is considered
satisfied.  For $q=\mu_M$, a term is finite only if $T\setminus i$ is
a basis of $M$.

Loops and parallel elements: loops are loops in every quotient.  If
$e,f$ are parallel in $M$, then the two-element circuit forces equality
of the corresponding coordinates in rank-one quotients; in higher rank
the flat-constancy lemma below forces equality for basis coordinates
differing only by replacing $e$ by $f$.  Thus the main question can be
studied on the simplification, then lifted by adding loop coordinates
and repeating parallel coordinates.

\subsection*{Structural consequence of the hypothesis}

A finite flat lattice is geometric and hence ranked and upper
semimodular.  If it has an anti-automorphism $x\mapsto x^\perp$, then
maximal chains are reversed, so $\rk(x^\perp)=r-\rk(x)$.  Applying the
anti-automorphism to the semimodular inequality gives the opposite
inequality; hence
\[
\rk(x)+\rk(y)=\rk(x\wedge y)+\rk(x\vee y)
\]
for all $x,y$.  Therefore $\cL(M)$ is modular.  In particular
\[
(x\vee y)^\perp=x^\perp\wedge y^\perp,\qquad
(x\wedge y)^\perp=x^\perp\vee y^\perp.
\]

\subsection*{Flat constancy of quotient coordinates}

If $\nu$ is a rank $k$ valuated quotient of $M$, then $\nu(B)$ depends
only on the flat $\cl_M(B)$ for independent $k$-sets $B$.

Proof: It suffices to compare adjacent bases
$A\cup\{b\}$ and $A\cup\{b'\}$ of the same rank $k$ flat $X$, where
$|A|=k-1$.  Choose $C$ of size $r-k$ such that
$A\cup\{b\}\cup C$ is a basis of $M$.  Since
$\cl(A b)=\cl(A b')=X$, also $A\cup\{b'\}\cup C$ is a basis.  Put
$T=A\cup\{b,b'\}\cup C$.  In
$\mathcal I_{k,r}(\nu,\mu_M;A,T)$ the $\mu_M$ factor is finite only
for $i=b,b'$; if $i\in C$, then $T-i$ contains the dependent set
$A\cup\{b,b'\}$ and has rank at most $r-1$.  The relation forces
$\nu(A b)=\nu(A b')$, including the $\infty$ cases.

Thus a rank $k$ quotient gives flat coordinates
\[
\bar\nu(X)=\nu(B)\quad\text{for }\rk X=k,\ \cl(B)=X,
\]
and $\nu(A)=\bar\nu(\cl A)$ for independent $A$ while $\nu(A)=\infty$
for dependent $A$.

\subsection*{Candidate orthogonal transform}

For a rank $k$ quotient $\nu$, define the candidate rank $r-k$
coordinate vector by
\[
\nu^{\perp_M}(A)=
\begin{cases}
\bar\nu((\cl A)^\perp),& A\text{ independent in }M,\ |A|=r-k,\\
\infty,&\text{otherwise}.
\end{cases}
\]
This corrects the dependent-subset issue: if $A$ is dependent, the
rank of $(\cl A)^\perp$ is larger than $k$, so there is in general no
rank $k$ flat coordinate to read.

The transform is involutive on flat coordinates.  For an embedded flat
$q_F=M/F\oplus U_{0,F}$, with $f=\rk F$, an independent $f$-set $A$
with $Y=\cl(A)$ is a basis after transforming exactly when
\[
Y^\perp\vee F=\hat1.
\]
Since $\rk(Y^\perp)=r-f$ and $\rk(F)=f$, modularity makes this
equivalent to $Y^\perp\wedge F=\hat0$, equivalently
$Y\vee F^\perp=\hat1$.  This is exactly the basis condition for
$q_{F^\perp}=M/F^\perp\oplus U_{0,F^\perp}$.

\subsection*{Main remaining lemma}

Need to prove the flat-coordinate duality lemma:

Let $L=\cL(M)$ be modular with opposition.  Let
$p_1,\ldots,p_s,\mu_M$ be a valuated flag of quotients of $M$, and
write each $p_j$ in flat coordinates.  Replacing each rank $k_j$
coordinate by $p_j^\perp(A)=\bar p_j((\cl A)^\perp)$ for independent
$A$ of size $r-k_j$ should produce a valuated flag of quotient ranks
$r-k_s<\cdots<r-k_1<r$.  Equivalently, all tropical Plucker and
incidence-Plucker relations for the flag should be invariant under
opposition after passing to flat coordinates.

Important correction: the flat-coordinate relations are not merely
``minimum twice on every rank-two interval''.  For $M=U_{2,4}$ and
rank-one quotients, the true condition is that the minimum is attained
at least twice on every three-element circuit.  The vector
$(0,0,1,1)$ has minimum twice on the unique rank-two flat $E$ but is
not in $\Trop(U_{2,4})$ because on the circuit $\{1,3,4\}$ the minimum
is unique.  Any proof must keep the subset/circuit-indexed relations.

Potential proof approaches:
\begin{itemize}
\item Relation-by-relation: express each Plucker/incidence relation in
terms of closures of the relevant subsets and show that the dual
collection of flat pairs is again a relation of the opposite ranks.
This likely uses modular interval duality and atomisticity.  Need an
explicit construction for arbitrary subset-indexed incidence relations.
\item Lattice cryptomorphism: Hirai's ``Uniform semimodular lattices
and valuated matroids'' gives a coordinate-free lattice-theoretic
description of valuated matroids and tropical linear spaces.  Need to
check whether it includes exactly the relative quotient/flag situation
and whether anti-automorphisms of a modular geometric lattice act on
the corresponding valuated objects.
\item Polar-space/building viewpoint: a finite modular geometric
lattice with opposition is a polar geometry; rank-$k$ quotient
valuations may be points in an affine building or valuated flag
geometry.  Opposition should be an isometry/anti-automorphism of this
building.  Need a precise theorem or a self-contained translation.
\item Classification: finite modular geometric lattices are products
of irreducible projective geometries (rank-two uniform lines, projective
planes, higher-rank Desarguesian geometries).  In coordinatizable
factors, Hodge star with respect to a polarity preserves the Plucker
and incidence ideals; tropicalization then preserves the Dressian.
Need to handle non-Desarguesian projective planes and products, and
also non-realizable valuated quotients.
\end{itemize}

\subsection*{Computational sanity checks}

Small checks support the candidate transform.

\begin{itemize}
\item $M=U_{2,n}$.  Rank-one quotients are points of the tropical line:
the minimum on every three-element circuit is attained at least twice.
A rank-two flat-lattice opposition permutes the point atoms, so the
candidate map is just the corresponding isometry of the tropical line
(on the rank-one layer), with top and bottom exchanged.
\item Fano plane $PG(2,2)$ with the standard dot-product polarity.
Random Bergman-fan points generated from flags were transformed to
rank-two flat-constant valuated matroids; direct checking of all
rank-two Plucker relations and all incidence relations with the Fano
matroid found no violations.
\item Direct sum $U_{2,3}\oplus U_{2,3}$ with componentwise rank-two
opposition.  For rank-one samples satisfying the two component circuit
conditions, the transformed rank-three coordinates passed all Plucker
and incidence checks.  For rank-two, all finite flat-coordinate values
in $\{0,1,2\}^{11}$ satisfying the quotient relations were enumerated;
the candidate opposition (which fixes the $3\times3$ cross-flats and
swaps the two component-top flats) preserved the solution set.
\end{itemize}

\subsection*{References to inspect}

Brandt--Eur--Zhang, \emph{Tropical flag varieties}, Adv. Math. 384
(2021), for valuated quotient/inclusion of projective tropical linear
spaces.

Dress--Wenzel, \emph{Valuated matroids}, Adv. Math. 93 (1992).

Hirai, \emph{Uniform semimodular lattices and valuated matroids},
JCTA 165 (2019), for a lattice-theoretic characterization of valuated
matroids and integer points in tropical linear spaces.  Need the exact
statements from the PDF/source.

Possible uniqueness references: Speyer, \emph{Tropical linear spaces};
Fink--Rincon or later tropical scheme papers for the statement that a
projective tropical linear space determines the valuated matroid up to
scalar via valuated cocircuits/covectors.

\end{document}
